\subsection{HMusic}\label{subsec:hmusic}

\textbf{HMusic} is a domain-specific language for music programming and live coding, embedded in the Haskell functional
programming language, and designed around the abstractions of patterns and tracks~\cite{Du_Bois_Ribeiro_2019}.

The language aims to bridge the gap between traditional music production tools, such as sequencers and drum machines,
and formal programming languages by providing a representation of music that resembles grid-based composition interfaces
while retaining the expressive power of functional programming~\cite{Du_Bois_Ribeiro_2019}.
Programs written in \textbf{HMusic} resemble the visual structure of multi-track sequencers, where patterns define
temporal structures and tracks associate these patterns with instruments, enabling both intuitive composition and
programmatic manipulation of musical structures~\cite{Du_Bois_Ribeiro_2019}.

%% ---------------------------------------------------------------------------------------------------------------------

\subsubsection{Key Features}

\textbf{HMusic} is centered around two main abstractions: \textit{patterns} and \textit{tracks}.
Patterns represent sequences of musical events, typically modeled as hits and rests, and are defined using
inductive data types that allow recursive composition.
Tracks associate patterns with instruments, enabling the separation of musical structure from sound generation.
Multiple tracks can be combined in parallel to form multi-track compositions, while sequential composition allows
the construction of longer musical structures from smaller components~\cite{Du_Bois_Ribeiro_2019}.

Example snippet:

\begin{lstlisting}[label={lst:hmusic-1}]
  kick = X :| O :| O :| O
  snare = O :| O :| X :| O
  track = MakeTrack "BassDrum" kick
    :|| MakeTrack "Snare" snare
\end{lstlisting}

As the language is embedded in Haskell, it benefits from functional programming features such as higher-order functions
and recursion, allowing users to define reusable transformations over patterns and tracks.
\textbf{HMusic} also supports live coding through primitives for playing, looping, and modifying tracks in real time,
enabling interactive and dynamic composition workflows.

Additionally, the language is compiled into \textit{Sonic Pi}, leveraging its synthesis capabilities while focusing
on higher-level musical abstractions~\cite{Du_Bois_Ribeiro_2019}.

%% ---------------------------------------------------------------------------------------------------------------------

\subsubsection{Limitations}

Despite its expressive abstractions, \textbf{HMusic} presents several limitations.

First, the language depends on \textit{Sonic Pi} as a backend for sound synthesis, limiting its independence and
requiring integration with external tools for execution.

Second, as an embedded DSL in Haskell, it inherits the complexity of functional programming, which may present a barrier
for users without prior programming experience.

Furthermore, the pattern-based model primarily focuses on rhythmic and grid-like structures, which may limit its ability
to represent more complex or hierarchical musical compositions.

Finally, while patterns and tracks provide compositional flexibility, the language does not explicitly model
higher-level temporal structures such as timelines or sections, which can make large-scale composition more difficult to
manage.

%% ---------------------------------------------------------------------------------------------------------------------

\subsubsection{Relevance to Our DSL}

\textbf{HMusic} demonstrates an effective approach to combining intuitive musical representations with formal
programming abstractions, particularly through its use of patterns and tracks as core compositional units.
Its design highlights the benefits of separating musical structure from sound generation and enabling composition
through reusable abstractions.

However, its reliance on pattern-based composition and external compilation to Sonic Pi contrasts with the goals of
this thesis, which aim to provide a more explicit and flexible temporal model.
The proposed DSL builds upon these ideas by introducing a timeline-based approach, improved modularity, and direct
compilation to formats such as MIDI or audio, enabling more structured and deterministic composition workflows.
